\section{Probabilidad}
\subsection{Axiomas}
\begin{enumerate}
    \item La probabilidad de un suceso debe de estar contenida en el rango $[0,1]$.
          \begin{equation*}
              0 \leq P(A)\leq 1
          \end{equation*}
    \item La probabilidad del conjunto vacio es 0.
          \begin{equation*}
              P(\emptyset)=0
          \end{equation*}
    \item La probabilidad del universo es 1.
          \begin{equation*}
              P(U)=1
          \end{equation*}
\end{enumerate}
\subsection{Eventos mutuamente excluyentes}
\begin{definicion}
    Sea A y B dos eventos, se dice que son mutuamente excluyentes si:
    \begin{equation*}
        A \cap B = \emptyset
    \end{equation*}
    $\cap$: Intersección, lo que quiere decir que A y B no tienen eventos en común, por lo tanto no pueden pasar al mismo tiempo.
\end{definicion}
\subsection{Addición de eventos}
La probabilidad que sucedan A y B es:
\begin{equation*}
    P(A\cup B) = P(A)+P(B) - P(A\cap B)
\end{equation*}
si A y B son eventos mutuamente excluyentes entonces:
\begin{align*}
    P(A\cup B) & = P(A)+P(B) - P(A\cap B) \\
               & =P(A)+P(B)-P(\emptyset)  \\
               & = P(A)+P(B)-0            \\
    P(A\cup B) & = P(A)+P(B)
\end{align*}
\subsection{Probabilidad condicional}
\begin{definicion}
    La probabilidad que suceda B dado que A ya sucedio, se lee como la probabilidad de B dado A y se denota como $P(B|A)$, y se calcula de la siguiente manera:
    \begin{equation*}
        P(B|A)= \frac{P(A\cap B)}{P(A)} \qquad \text{si} \qquad P(A)>0
    \end{equation*}
\end{definicion}
Si A y B son dos eventos independientes entonces:
\begin{align*}
    P(B|A) & = \frac{P(A\cap B)}{P(A)} \\
           & = \frac{P(A)P(B)}{P(A)}   \\
           & = P(B)
\end{align*}
\begin{ejemplo}
    La probabilidad de que un vuelo programado en forma regular salga a tiempo es $P(D)=0.83$, la probabilidad de que llegue a tiempo es $P(A)=0.92$ y la probabilidad de que salga y llegue a tiempo es $P(D\cup A)=0.78$. Encuentre la probabilidad de que un avión:
    \begin{itemize}
        \item Llegue a tiempo dado que salió a tiempo.
              \begin{align*}
                  P(A|D) & = \frac{P(A\cap D)}{P(D)} \\
                         & = \frac{0.78}{0.83}       \\
                         & = 0.94                    \\
                         & = 94\%
              \end{align*}
        \item Haya salido a tiempo dado que llego a tiempo.
              \begin{align*}
                  P(D|A) & = \frac{P(A\cap D)}{P(D)} \\
                         & = \frac{0.78}{0.92}       \\
                         & = 0.847                   \\
                         & = 84.7\%
              \end{align*}
    \end{itemize}
\end{ejemplo}
\begin{ejemplo}
    Encuentre la probabilidad de elegir al azar y en forma sucesiva 4 litros buenos de leche de un refrigerador que contiene 20 litros de los cuales 5 litros están en mal estado.Sean: \begin{itemize}
        \item $A_1$: El primer litro este en buen estado.
        \item $A_2$: El segundo litro este en buen estado.
        \item $A_3$: El tercero litro este en buen estado.
        \item $A_4$: El cuarto litro este en buen estado.
    \end{itemize}
    Entonces:
    \begin{align*}
        P(A_1 \cup A_2 \cup A_3 \cup A_4) & = P(A_1)P(A_2|A_1)P(A_3|A_2,A_1)P(A_4|A_3,A_2,A_1)                                                         \\
                                          & = \left(\frac{15}{20}\right)\left(\frac{14}{19}\right)\left(\frac{13}{18}\right)\left(\frac{12}{17}\right) \\
                                          & = \frac{91}{323}
    \end{align*}
\end{ejemplo}
\subsection{Probabilidad total para varios eventos}
\begin{definicion}
    Sean $A_1,A_2,\dots,A_k$ eventos, se dice que son exhaustivos si
    \begin{equation*}
        A_1 \cup A_2 \cup \dots \cup A_k =U
    \end{equation*}
\end{definicion}
Sean $A_1,A_2,\dots,A_k$ eventos mutuamente excluyentes y exhaustivos, entonces
\begin{equation*}
    P(B)=P(B\cup A_1) P(B\cup A_2)+P(B\cup)+\dots+P(B\cup A_k)
\end{equation*}
\begin{ejemplo}
    Una empresa consultora renta automóviles de tres agencias, 20\% de la agencia D, 20\% de la agencia E y 60\% de la agencia F. Si 10\% de los autos de D, 12\% de los autos de E y 4\% de los autos de F tienen neumáticos en mal estado. ¿Cuál es la probabilidad de que la empresa reciba autos con neumáticos en mal estado?
    Sean los eventos:
    \begin{itemize}
        \item $A_1$: Automóvil rentado proviene de la agencia D
        \item $A_2$: Automóvil rentado proviene de la agencia E
        \item $A_3$: Automóvil rentado proviene de la agencia F
        \item B: Automóvil rentado tiene neumáticos en mal estado.
    \end{itemize}
    \begin{align*}
        P(B) & =P(B\cup A_1)+P(B\cup A_2)+P(B\cup A_3)       \\
             & =P(A_1)P(B|A_1)+P(A_2)P(B|A_2)+P(A_3)P(B|A_3) \\
             & =(0.2)(0.1)+(0.2)(0.12)+(0.6)(0.04)           \\
             & =0.068
    \end{align*}
\end{ejemplo}
\begin{ejemplo}
    Supongamos que 5\% de todos los hombres y 0.25\% de todas las mujeres son daltónicas. Una persona elegida al azar resulta ser daltónica. ¿Cuál es la probabilidad de que esta persona sea hombre? (se considera que la cantidad de hombres y mujeres es igual). Sean los eventos:
    \begin{itemize}
        \item $A_1$: La persona elegida sea un hombre
        \item $A_2$: La persona elegida sea una mujer
        \item B: La persona elegida es daltónica.
    \end{itemize}
    \begin{table}[H]
        \centering
        \begin{tabular}{lcccc}
            \hline
            Eventos & $P(A_i)$ & $P(B|A_i)$ & $P(B\cup A_i)$ & $P(A_i|B)$ \\ \hline
            $A_1$   & 0.5      & 0.05       &                &            \\
            $A_2$   & 0.5      & 0.0025     &                &            \\
            Total   &          & 0.02625    &                             \\ \hline
        \end{tabular}
        \caption{Información del ejemplo}
    \end{table}
    Como se pregunta $P(A_1|B)$, entonces tendremos que calcular lo siguiente:
    \begin{equation*}
        P(A_1|B)= \frac{P(A_1 \cup B)}{P(B)}
    \end{equation*}
    como no tenemos cuanto es $P(B)$ lo tendremos que calcular de la siguiente manera:
    \begin{equation*}
        P(B)=P(A_1 \cup B) + P(A_2 \cup B)
    \end{equation*}
    calculando $P(A_1 \cup B)$ se obtiene que:
    \begin{equation*}
        P(A_1 \cup B) = P(A_1)P(B|A_1) =(0.5)(0.05)= 0.025
    \end{equation*}
    calculando $P(A_2 \cup B)$ se obtiene que:
    \begin{equation*}
        P(A_2 \cup B) = P(A_2)P(B|A_2) =(0.5)(0.025) =0.00125
    \end{equation*}
    entonces:
    \begin{align*}
        P(B) & =P(A_1)P(A_1 \cup B) + P(A_2)P(A_2 \cup B) \\
             & = (0.025)+(0.00125)                        \\
             & = 0.02625
    \end{align*}
    por lo tanto:
    \begin{align*}
        P(A_1|B) & = \frac{P(A_1 \cup B)}{P(B)} \\
                 & = \frac{0.025}{0.02625}      \\
                 & = 0.9523
    \end{align*}
\end{ejemplo}
\subsection{Teorema de Bayes}
El teorema de Bayes es una rama de la probabilidad condicional, nos permite conocer la probabilidad de una causa a partir de la probabilidad de la consecuencia.
\begin{ejemplo}
    En cierto sector metropolitano los grupos de ingresos bajos, medios y altos, constituyen el 20\%, 55\% y 25\% de la población respectivamente. Se sabe además que el 80\% del grupo de bajos ingresos se opone a un proyecto de ley que cursa actualmente en el congreso. Los porcentajes de los grupos medios y alto son 30 y 10\% respectivamente. Se selecciona al azar un individuo entre la población y se encuentra que se opone al proyecto de ley. ¿Cuál es la probabilidad de que la persona pertenezca al grupo de bajos ingresos? ¿Al de ingresos medios? ¿Al de altos ingresos?
    Sean los eventos:
    \begin{itemize}
        \item $A_1$: Que la persona seleccionada pertenezca al grupo de ingresos bajos.
        \item $A_2$: Que la persona seleccionada pertenezca al grupo de ingresos medios.
        \item $A_3$: Que la persona seleccionada pertenezca al grupo de ingresos altos.
        \item B: Que la persona elegida se opone al proyecto de ley.
    \end{itemize}
    \begin{table}[H]
        \centering
        \begin{tabular}{lcc}\hline
            Eventos & $P(A_i)$ & $P(B|A_i)$ \\ \hline
            $A_1$   & 0.2      & 0.8        \\
            $A_2$   & 0.55     & 0.3        \\
            $A_3$   & 0.25     & 0.1        \\\hline
        \end{tabular}
        \caption{Información del ejemplo}
    \end{table}
    Lo que nos preguntan en la probabilidad de que sea del grupo $A_i$ dado que se opone a la propuesta, eso sería $P(A_i|B)$, entonces tendriamos que calcular:
    \begin{equation*}
        P(A_i|B)=\frac{P(A_i \cup B)}{P(B)}
    \end{equation*}
    calculando $P(B)$ se tiene lo siguiente:
    \begin{equation*}
        P(B) = P(B\cup A_1)+P(B\cup A_2)+P(B\cup A_3)
    \end{equation*}
    como no conocemos cuando valen $P(B\cup A_i)$, las calcularemos primero, eso seria:
    \begin{align*}
        P(B\cup A_1) & = P(A_1)P(B|A_1) = (0.2)(0.8)=0.160  \\
        P(B\cup A_2) & = P(A_2)P(B|A_2) = (0.55)(0.3)=0.165 \\
        P(B\cup A_3) & = P(A_3)P(B|A_3) = (0.25)(0.1)=0.025
    \end{align*}
    entonces:
    \begin{align*}
        P(B) & = P(B\cup A_1)+P(B\cup A_2)+P(B\cup A_3) \\
             & = 0.160+0.165+0.025                      \\
             & = 0.350
    \end{align*}
    por lo tanto:\\
    \begin{minipage}{0.3\linewidth}
        \begin{align*}
            P(A_1|B) & =\frac{P(A_1 \cup B)}{P(B)} \\
                     & = \frac{0.160}{0.350}       \\
                     & = 0.4571
        \end{align*}
    \end{minipage}
    \begin{minipage}{0.3\linewidth}
        \begin{align*}
            P(A_2|B) & =\frac{P(A_2 \cup B)}{P(B)} \\
                     & = \frac{0.165}{0.350}       \\
                     & = 0.4714
        \end{align*}
    \end{minipage}
    \begin{minipage}{0.3\linewidth}
        \begin{align*}
            P(A_3|B) & =\frac{P(A_3 \cup B)}{P(B)} \\
                     & = \frac{0.025}{0.350}       \\
                     & = 0.0714
        \end{align*}
    \end{minipage}
\end{ejemplo}